\documentclass[11pt]{article}
\usepackage[margin=25mm]{geometry}
\usepackage{fontspec}
\setmainfont{Latin Modern Roman}
\usepackage{amsmath,amssymb,mathtools}
\usepackage{unicode-math}
\setmathfont{Latin Modern Math}
\usepackage{xurl,hyperref}
\hypersetup{colorlinks=true,urlcolor=blue}
\setcounter{secnumdepth}{0}
\setlength{\parindent}{0pt}
\setlength{\parskip}{5pt}
\setlength{\emergencystretch}{3em}
\newcommand{\sref}[2]{\hyperlink{source:#1:#2}{[S#2]}}
\newcommand{\cataloguescope}[1]{\paragraph{Recorded target.}#1}
\begin{document}
\section{Multiple ergodic averages — Problem 30}\label{statement}

\subsection{Definitions and mathematical statement}

\paragraph{Shared notation.}$\mathbb N=\{1,2,\ldots\}$, and $\mathbb Z,\mathbb Q,\mathbb R,\mathbb C$ denote the integers, rationals, reals and complex numbers. Any other conventions are specified locally.


Let $\mathbb P$ be the positive prime numbers, $\ell\in\mathbb N$, and $c,c_1,\ldots,c_\ell>0$ be arbitrary real numbers. Floors $\lfloor t\rfloor$ are integer parts; the exponents need not be distinct or nonintegral.
\paragraph{Statement recorded in the supplied source.}
Do the primes contain both types of patterns \[\{m,m+\lfloor n^c\rfloor,m+2\lfloor n^c\rfloor,\ldots,m+\ell\lfloor n^c\rfloor\},\qquad \{m,m+\lfloor n^{c_1}\rfloor,\ldots,m+\lfloor n^{c_\ell}\rfloor\}?\] Each type is required to occur for infinitely many $n\in\mathbb N$, with a suitable starting value $m\in\mathbb N$ allowed to depend on $n$ and the type. \sref{statement}{2}
\subsection{Short English statement}
Do primes contain arbitrarily long progressions with power-floor differences and all specified mixed-power-floor patterns, for infinitely many input integers?
\subsection{Sources}
\begin{description}
\item[\hypertarget{source:statement:1}{S1}] ULAM.AI and the UnsolvedMath Contributors, \emph{UnsolvedMath: A Curated Collection of Open Mathematics Problems}, catalogue entry. CC BY 4.0. \url{https://huggingface.co/datasets/ulamai/UnsolvedMath}.
\item[\hypertarget{source:statement:2}{S2}] N. Frantzikinakis, Some open problems on multiple ergodic averages (2016), Section 6.3, full Problem 30, p.34. \url{https://arxiv.org/pdf/1103.3808v3}.
\end{description}
\label{end:statement}
\end{document}
